% Bài làm buổi 11/11
\begin{problem}[9DS-HSB2-013]
a) Thay $t = 3$: $s = 5 \cdot 3^2 = 5 \cdot 9 = 45$ (m).

b) Gọi $t$ (giây, $t > 0$) là thời gian vật rơi đến khi chạm đất. Khi đó $5t^2 = 125$.

$t^2 = 25$, suy ra $t = 5$ hoặc $t = -5$. Vì $t > 0$ nên $t = 5$ (giây).

c) Sau $2$ giây vật rơi được $s = 5 \cdot 2^2 = 20$ (m), vật chưa chạm đất vì $2 < 5$.

Vật còn cách mặt đất $125 - 20 = 105$ (m).

Đáp số: a) $45$ m; b) $5$ giây; c) $105$ m
\end{problem}

\begin{problem}[9DS-HSB2-001]
a) $f(-3) = 2 \cdot (-3)^2 = 2 \cdot (-9) = -18$.

$f\left(\dfrac{1}{2}\right) = 2 \cdot \left(\dfrac{1}{2}\right)^2 = 2 \cdot \dfrac{1}{4} = \dfrac{1}{2}$.

$f\left(\sqrt{2}\right) = 2 \cdot \left(\sqrt{2}\right)^2 = 2 \cdot 2 = 4$.

b) $y = 8$ nên $2x^2 = 8$, tức là $x^2 = 4$.

Suy ra $x = 2$ hoặc $x = -2$.

Đáp số: a) $f(-3) = -18$, $f\left(\dfrac{1}{2}\right) = \dfrac{1}{2}$, $f\left(\sqrt{2}\right) = 4$; b) $x = 2$; $x = -2$
\end{problem}
